\documentclass[a4paper, 14pt]{extarticle}

% Поля
%--------------------------------------
\usepackage{geometry}
\geometry{a4paper,tmargin=2cm,bmargin=2cm,lmargin=3cm,rmargin=1cm}
%--------------------------------------

\usepackage{float}
%Russian-specific packages
%--------------------------------------
\usepackage[T2A]{fontenc}
\usepackage[utf8]{inputenc} 
\usepackage[english, main=russian]{babel}
%--------------------------------------

\usepackage{textcomp}

% Красная строка
%--------------------------------------
\usepackage{indentfirst}               
%--------------------------------------             
\usepackage[most]{tcolorbox}
\tcbuselibrary{listings,breakable}
\newtcblisting{breakablelisting}{
  listing only,
  breakable,
  enhanced,
  frame hidden,
  colback=white,
  listing options={
    mathescape=false,
    texcl=false,
    inputencoding=utf8,
    literate={\$}{{\$}}1
  }
}


%Graphics
%--------------------------------------
\usepackage{graphicx}
\usepackage{wrapfig}
%--------------------------------------
\graphicspath{{./lab5/}}
% Полуторный интервал
%--------------------------------------
\linespread{1.3}                    
%--------------------------------------

%Выравнивание и переносы
%--------------------------------------
% Избавляемся от переполнений
\sloppy
% Запрещаем разрыв страницы после первой строки абзаца
\clubpenalty=10000
% Запрещаем разрыв страницы после последней строки абзаца
\widowpenalty=10000
%--------------------------------------

%Списки
\usepackage{enumitem}

%Подписи
\usepackage{caption} 

%Гиперссылки
\usepackage{hyperref}

\hypersetup {
	unicode=true
}

%Рисунки
%--------------------------------------
\DeclareCaptionLabelSeparator*{emdash}{~--- }
\captionsetup[figure]{labelsep=emdash,font=onehalfspacing,position=bottom}
%--------------------------------------

\usepackage{tempora}

%Листинги
%--------------------------------------
\usepackage{listings}
\lstset{
  basicstyle=\ttfamily\footnotesize, 
  %basicstyle=\footnotesize\AnkaCoder,        % the size of the fonts that are used for the code
  breakatwhitespace=false,         % sets if automatic breaks shoulbd only happen at whitespace
  breaklines=true,                 % sets automatic line breaking
  captionpos=t,                    % sets the caption-position to bottom
  inputencoding=utf8,
  frame=single,                    % adds a frame around the code
  keepspaces=true,                 % keeps spaces in text, useful for keeping indentation of code (possibly needs columns=flexible)
  keywordstyle=\bf,       % keyword style
  numbers=left,                    % where to put the line-numbers; possible values are (none, left, right)
  numbersep=5pt,                   % how far the line-numbers are from the code
  xleftmargin=25pt,
  xrightmargin=25pt,
  showspaces=false,                % show spaces everywhere adding particular underscores; it overrides 'showstringspaces'
  showstringspaces=false,          % underline spaces within strings only
  showtabs=false,                  % show tabs within strings adding particular underscores
  stepnumber=1,                    % the step between two line-numbers. If it's 1, each line will be numbered
  tabsize=2,                       % sets default tabsize to 8 spaces
  title=\lstname                   % show the filename of files included with \lstinputlisting; also try caption instead of title
}
%--------------------------------------

%%% Математические пакеты %%%
%--------------------------------------
\usepackage{amsthm,amsfonts,amsmath,amssymb,amscd}  % Математические дополнения от AMS
\usepackage{mathtools}                              % Добавляет окружение multlined
\usepackage[perpage]{footmisc}
%--------------------------------------

%--------------------------------------
%			НАЧАЛО ДОКУМЕНТА
%--------------------------------------

\begin{document}

%--------------------------------------
%			ТИТУЛЬНЫЙ ЛИСТ
%--------------------------------------
\begin{titlepage}
\thispagestyle{empty}
\newpage


%Шапка титульного листа
%--------------------------------------
\vspace*{-60pt}
\hspace{-65pt}
\begin{minipage}{0.3\textwidth}
\hspace*{-20pt}\centering
\includegraphics[width=\textwidth]{emblem}
\end{minipage}
\begin{minipage}{0.67\textwidth}\small \textbf{
\vspace*{-0.7ex}
\hspace*{-6pt}\centerline{Министерство науки и высшего образования Российской Федерации}
\vspace*{-0.7ex}
\centerline{Федеральное государственное автономное образовательное учреждение }
\vspace*{-0.7ex}
\centerline{высшего образования}
\vspace*{-0.7ex}
\centerline{<<Московский государственный технический университет}
\vspace*{-0.7ex}
\centerline{имени Н.Э. Баумана}
\vspace*{-0.7ex}
\centerline{(национальный исследовательский университет)>>}
\vspace*{-0.7ex}
\centerline{(МГТУ им. Н.Э. Баумана)}}
\end{minipage}
%--------------------------------------

%Полосы
%--------------------------------------
\vspace{-25pt}
\hspace{-35pt}\rule{\textwidth}{2.3pt}

\vspace*{-20.3pt}
\hspace{-35pt}\rule{\textwidth}{0.4pt}
%--------------------------------------

\vspace{1.5ex}
\hspace{-35pt} \noindent \small ФАКУЛЬТЕТ\hspace{80pt} <<Информатика и системы управления>>

\vspace*{-16pt}
\hspace{47pt}\rule{0.83\textwidth}{0.4pt}

\vspace{0.5ex}
\hspace{-35pt} \noindent \small КАФЕДРА\hspace{50pt} <<Теоретическая информатика и компьютерные технологии>>

\vspace*{-16pt}
\hspace{30pt}\rule{0.866\textwidth}{0.4pt}
  
\vspace{11em}

\begin{center}
\Large {\bf Летучка № 2} \\ 
\large {\bf по курсу <<Методы оптимизации>>} \\
\large <<Метод простых симплексов>> 
\end{center}\normalsize

\vspace{8em}


\begin{flushright}
  {Студент группы ИУ9-82Б Шемякин В.А. \hspace*{15pt}\\ 
  \vspace{2ex}
  Преподаватель Посевин Д. П.\hspace*{15pt}}
\end{flushright}

\bigskip

\vfill
 

\begin{center}
\textsl{Москва 2026}
\end{center}
\end{titlepage}
%--------------------------------------
%		КОНЕЦ ТИТУЛЬНОГО ЛИСТА
%--------------------------------------

\renewcommand{\ttdefault}{pcr}

\setlength{\tabcolsep}{3pt}
\newpage
\setcounter{page}{2}

\section{Задание}\label{Sect::task}
\begin{enumerate}
\item Реализовать простой симплекс метод
\item Протестировать его на примере функций Швеффеля, Растригина и Розенброкка.
\end{enumerate}

\section{Результаты}\label{Sect::res}

Исходный код программы представлен в листинге 1.

\begin{breakablelisting}
    using LinearAlgebra
    using Plots
    Plots.default(background=:white, foreground_color=:darkgray)
    
    function initial_simplex1(n)
        P = (1 / (n*sqrt(2))) * (sqrt(n + 1) + n - 1)
        Q = (1 / (n*sqrt(2))) * (sqrt(n + 1) - 1)
        X = zeros(n, n + 1)
        for k in 2:n+1
            for j in 1:n
                X[j, k] = (j == k - 1) ? P : Q
            end
        end
        return X
    end
    
    function initial_simplex2(n)
        R = [sqrt(1 / (2 * i*(i + 1))) for i in 1:n]
        V = [sqrt(i / (2 * (i + 2))) for i in 1:n]
        X = zeros(n, n + 1)
        for j in 1:n
            X[j, 1] = -R[j]
        end
        for k in 2:n+1
            for j in 1:n
                X[j, k] = (j == k - 1) ? V[k-1] : -R[j]
            end
        end
        return X
    end
    
    function simple_simplex(f, x0; ε=1e-6, maxiter=10000, σ=0.5, history=nothing, init_method=1, ref_max_ratio=2.0)
        n = length(x0)
        S = (init_method == 1) ? initial_simplex1(n) : initial_simplex2(n)
        X = x0 .+ S
        fvals = [f(X[:, i]) for i in 1:n+1]
        iter = 0
        if history !== nothing
            push!(history, copy(X))
        end
        while iter < maxiter
            order = sortperm(fvals)
            fvals = fvals[order]
            X = X[:, order]
            x_centroid = sum(X[:, 1:n], dims=2) / n
            x_ref = 2 * x_centroid - X[:, n+1]
            d_worst = norm(X[:, n+1] - x_centroid)
            d_ref = norm(x_ref - x_centroid)
            if d_worst > 1e-12 && d_ref > ref_max_ratio * d_worst
                x_ref = x_centroid + (x_ref - x_centroid) .* (ref_max_ratio * d_worst / d_ref)
            end
            f_ref = f(x_ref)
            if f_ref < fvals[n+1]
                X[:, n+1] = x_ref
                fvals[n+1] = f_ref
            else
                for i in 2:n+1
                    X[:, i] = X[:, 1] + σ * (X[:, i] - X[:, 1])
                    fvals[i] = f(X[:, i])
                end
            end
            if history !== nothing
                push!(history, copy(X))
            end
            if sqrt(sum(sum((X[:, i] .- x_centroid).^2) for i in 1:n+1) / (n+1)) < ε
                break
            end
            iter += 1
        end
        return X[:, 1], fvals[1], iter
    end
    
    schwefel(x) = 418.9829 * length(x) - sum(x .* sin.(sqrt.(abs.(x))))
    rastrigin(x) = 10 * length(x) + sum(x.^2 .- 10 .* cos.(2π .* x))
    rosenbrock(x) = sum(100 * (x[i+1] - x[i]^2)^2 + (1 - x[i])^2 for i in 1:length(x)-1)
    
    function plot_simplex_history(f, history, xlims, ylims, title_str, filename)
        all_x = Float64[]
        all_y = Float64[]
        for X in history
            append!(all_x, X[1, :])
            append!(all_y, X[2, :])
        end
        margin_x = max(0.1 * (maximum(all_x) - minimum(all_x)), 0.5)
        margin_y = max(0.1 * (maximum(all_y) - minimum(all_y)), 0.5)
        xmin = minimum(all_x) - margin_x
        xmax = maximum(all_x) + margin_x
        ymin = minimum(all_y) - margin_y
        ymax = maximum(all_y) + margin_y
        xmin = max(xmin, xlims[1])
        xmax = min(xmax, xlims[2])
        ymin = max(ymin, ylims[1])
        ymax = min(ymax, ylims[2])
        p = plot(background=:white, xlims=(xmin, xmax), ylims=(ymin, ymax))
        for (k, X) in enumerate(history)
            tri_x = [X[1, 1]; X[1, 2]; X[1, 3]; X[1, 1]]
            tri_y = [X[2, 1]; X[2, 2]; X[2, 3]; X[2, 1]]
            plot!(p, tri_x, tri_y, linewidth=2.5, color=:black, alpha=0.85, legend=false)
            if k == 1
                scatter!(p, X[1, :], X[2, :], markersize=4, color=:black, marker=:circle, legend=false)
            end
        end
        if length(history) >= 1
            X1 = history[1]
            annotate!(p, [(X1[1, j], X1[2, j], text(string(j), 8, :center)) for j in 1:3])
            scatter!(p, [X1[1, 1]], [X1[2, 1]], markersize=4, color=:blue, label="начало")
        end
        X = history[end]
        scatter!(p, [X[1, 1]], [X[2, 1]], markersize=4, color=:red, label="минимум")
        plot!(p, xlims=(xmin, xmax), ylims=(ymin, ymax), title=title_str, xlabel="x₁", ylabel="x₂")
        savefig(p, filename)
    end
    
    n = 2
    x0_s = [400.0, 400.0]
    x0_r = [3.0, 3.0]
    x0_ros = [-1.5, 2.0]
    
    println("Швеффель n=$n")
    hist_s = []
    x, f, it = simple_simplex(schwefel, x0_s, history=hist_s)
    println("Точка: $x")
    println("Значение функции: $f")
    println("Итераций: $it")
    
    println("\nРастригин n=$n")
    hist_r = []
    x, f, it = simple_simplex(rastrigin, x0_r, history=hist_r)
    println("Точка: $x")
    println("Значение функции: $f")
    println("Итераций: $it")
    
    println("\nРозенброк n=$n")
    hist_ros = []
    x, f, it = simple_simplex(rosenbrock, x0_ros, history=hist_ros)
    println("Точка: $x")
    println("Значение функции: $f")
    println("Итераций: $it")
    
    plot_simplex_history(schwefel, hist_s, (-500, 500), (-500, 500), "Швеффель", "let2_schwefel.png")
    plot_simplex_history(rastrigin, hist_r, (-5, 5), (-5, 5), "Растригин", "let2_rastrigin.png")
    plot_simplex_history(rosenbrock, hist_ros, (-2, 2.5), (-1, 2.5), "Розенброк", "let2_rosenbrock.png")
    println("\nГрафики сохранены: let2_schwefel.png, let2_rastrigin.png, let2_rosenbrock.png")
    
\end{breakablelisting}

\newpage

Результат представлен на рисунках~\ref{fig:img1}--~\ref{fig:rastrigin}.

\begin{figure}[H]
  \centering
  \includegraphics[width=0.6\textwidth]{let2/screen1.png}
  \caption{Результат}
  \label{fig:img1}
\end{figure}

\begin{figure}[H]
  \centering
  \includegraphics[width=0.6\textwidth]{let2/screen2.png}
  \caption{Результат}
  \label{fig:schwefel}
\end{figure}

\begin{figure}[H]
  \centering
  \includegraphics[width=0.6\textwidth]{let2/screen3.png}
  \caption{Результат}
  \label{fig:rastrigin}
\end{figure}

\end{document}
